% Front matter for the continuously paginated retrospective.
% The binder builder supplies the paper counts, page ranges and contents rows.
\documentclass[11pt,letterpaper]{article}
\usepackage[T1]{fontenc}
\usepackage[utf8]{inputenc}
\usepackage{lmodern,microtype,booktabs,array,xcolor,fancyhdr}
\usepackage[letterpaper,margin=1in,includeheadfoot,headheight=12pt,
  headsep=12pt,footskip=20pt]{geometry}
\usepackage[hidelinks]{hyperref}
\microtypesetup{protrusion=false}
\hypersetup{pdftitle={Curious Signals: Measurements and Laboratory Models},
  pdfauthor={Sebastian J. Spicker, Alexander K\"upper},
  pdfsubject={A retrospective of work from 2018 to 2024; October 2026}}
\setlength{\parindent}{0pt}
\setlength{\parskip}{0.62em}
\setlength{\emergencystretch}{2em}
\renewcommand{\headrulewidth}{0pt}
\renewcommand{\footrulewidth}{0pt}
\fancyhf{}
\fancyhead[L]{\footnotesize\textsc{Curious Signals}}
\fancyhead[R]{\footnotesize\thepage}
\pagestyle{fancy}
\newcommand{\paperlink}[2]{\href{cs-paper:#1}{#2}}
\newcommand{\paperpages}[1]{\paperlink{#1}{\csname paperpages@#1\endcsname}}
\newcommand{\contentsgroup}[1]{\normalsize\bfseries #1}
\newcommand{\overviewheading}[1]{%
  \par\vspace{0.65em}{\large\bfseries #1\par}\vspace{0.12em}}
\newcommand{\fronttitle}[1]{%
  {\fontsize{28}{33}\selectfont #1\par}\vspace{0.55em}}
% PAPER_PAGE_DEFINITIONS
\begin{document}
\thispagestyle{empty}
{\small\sffamily INSTITUTE OF PHYSICS EDUCATION\\UNIVERSITY OF COLOGNE}

\vspace*{1.0in}
{\fontsize{44}{49}\selectfont Curious Signals\par}
\vspace{0.23in}
{\fontsize{25}{31}\selectfont Measurements and\\laboratory models\par}

\vspace{0.5in}
{\large A retrospective of work from 2018 to 2024\par}

\vspace{0.48in}
{\large Sebastian J. Spicker\\[0.15em]Alexander K\"upper\par}

\vfill
{\normalsize
% PAPER_COUNT
papers on sensing, measurement and astronomical inference\par}
\vspace{0.22in}
{\small October 2026}

\newpage
\fancyhead[L]{\footnotesize\textsc{Retrospective}}
\fronttitle{From signals to physical quantities}

Curious Signals connects classroom sensor experiments with the physical models
used to interpret them. Its twenty-one papers examine how calibration, timing and
geometry relate recorded signals to physical quantities. Each paper makes that
relationship explicit through a measurement model, a reference comparison and a
worked reduction.

The papers cover forty-three experiment definitions and analyse fifty-eight
records. The reported results connect sensor readings to physical quantities
through explicit conversion, calibration and timing conventions. Apparatus
schematics identify the reference geometry, and each record retains its own time
coordinates. The analyses distinguish observed variation from systematic effects
and state the assumptions used to infer quantities beyond the recorded signals.

\overviewheading{Motion, fields and electrical measurements}
The inertial studies use vector information to resolve orientation and motion.
\paperlink{accelerometer}{Specific-force direction} gives static inclination
when translational acceleration is negligible, while magnitude and laboratory
alignment require separate references. \paperlink{gyroscope}{Gyroscope rate
reversals} identify scale and offset only against independent reference rates;
signed integrals distinguish net rotation about a fixed axis from accumulated
motion, and a residual bias grows into an angle error over time. Magnetic
calibration, heading and \paperlink{magnetometer}{indirect current measurement}
each require their own reference information and stable geometry.

The electrical studies separate numerical conversion from calibration.
\paperlink{analog-inputs}{ADC codes, voltage plateaux and RC curves} expose
different roles for granularity, load and the timebase. A friction estimate
depends on the \paperlink{static-friction}{inclination before first slip}, while
contact history and mounting enter the inferred coefficient.
\paperlink{audio-magnetic}{Sound envelopes and magnetic fields} track changes
in a motor's operating state through complementary responses.

\overviewheading{Pressure, humidity and thermal response}
Pressure connects to height, force and gas state through different boundary
conditions. \paperlink{pressure}{Bracketed ambient references} resolve small
height-related changes, while \paperlink{mission-to-mars}{enclosure and
compression examples} require temperature, volume or area information to close
the model; a slow pressure relaxation, for example, may reflect gas exchange or
cooling.

\paperlink{temperature-humidity}{Paired temperature and humidity channels}
distinguish warming from a change in water-vapour content. Enclosure-heating
curves quantify contrasts and apparent response times, but a fitted response time
does not identify a \paperlink{greenhouse}{heat-loss coefficient}; separating
optical from convective effects requires controlled source power and independent
air thermometry. The \paperlink{pressure-boiling}{boiling study} similarly
separates a visible observational threshold from equilibrium saturation, while
accounting for instrument range and the applicable thermodynamic domain.

\newpage
\fronttitle{Light, geometry and inference}

\overviewheading{Optics and radiometry}
Optical ratios become interpretable once the reference, detector weighting and
viewing geometry are specified. The \paperlink{light-rgb}{RGB, LED and UV
examples} separate detector-weighted transmission from an absolute radiometric
scale. \paperlink{illuminance-distance}{Distance studies} recover exponents close
to two and quantify the roles of averaging and the source-origin convention.
\paperlink{albedo}{Reflected-light ratios} depend on the illumination and viewing
geometry, whereas Bond albedo integrates reflected power over angle and stellar
spectrum.

\paperlink{malus-law}{Settled-angle polarimetry} separates scale, floor and
angular zero in a cosine-squared response, with the angular reference entering
the interpretation of the fitted phase. The
\paperlink{colour-temperature}{colour-temperature study} obtains a consistent
colour ordering through a transferred sensor matrix, but its temperature scale
inherits that matrix's calibration. \paperlink{ir-distance}{Infrared radiometry}
relates temperature to emissivity, scene coverage and detector thermal state;
outward and return branches reveal the role of acquisition history. In
\paperlink{spectral-fingerprints}{paired spectra}, removal of a shared feature
separates total attenuation from additional attenuation. Wavelength calibration,
spectral resolution and detection significance enter as distinct quantities.

\overviewheading{Planetary and astronomical models}
The astronomical studies retain the same order of inference: constrain the
forcing and detector response, then interpret the physical model.
\paperlink{tidal-locking}{Illumination and thermal-contrast examples} separate
incidence geometry from absorbed power, heat redistribution and sensor
response.

\paperlink{transit-method}{Laboratory transit photometry} brings together
detectors, actual timestamps and visible-area models. Local reference windows
recover shallow deficits near four percent in the worked examples, while
whole-record extrema answer a different question. Coarse sampling can miss a
short feature altogether. Planet--moon overlaps, dark source patches and pairs
of luminous bodies alter the brightness integral in different ways, so no single
conversion turns a light-curve minimum into an object radius.
The \paperlink{coronagraph-contrast}{occulting-mask study} completes the theme:
background suppression, target throughput and noise together determine
detectability.

\overviewheading{A common experimental discipline}
The collection treats calibration, geometry, time response and uncertainty as
parts of the physical argument. Sample variability, shared offsets and correlated
errors enter the uncertainty of a derived quantity through its measurement model;
scatter within a record does not bound its systematic error. Across the studies,
the reported comparisons connect the observed signals to the inferred quantities
and their uncertainties. Independent references and explicit model assumptions
establish the scope of each physical interpretation.

\vfill
{\footnotesize
Each paper retains its own references and date. Page numbers run continuously
through this volume; the contents and PDF bookmarks locate every article.
The core experiment software was developed by Gautier Creutzer and
Fr\'ed\'eric Bouquet, with revisions by Sebastian J. Spicker and
Fr\'ed\'eric Bouquet.}

\newpage
\fancyhead[L]{\footnotesize\textsc{Contents}}
\fronttitle{Contents}
{\small
% PAPER_COUNT
papers\quad\textperiodcentered\quad
% BODY_PAGE_COUNT
article pages\par}
\vspace{1.0em}
\begingroup
\small
\renewcommand{\arraystretch}{1.2}
\begin{tabular}{@{}p{0.68\textwidth}@{\hspace{0.03\textwidth}}p{0.18\textwidth}@{\hspace{0.03\textwidth}}>{\raggedleft\arraybackslash}p{0.08\textwidth}@{}}
\toprule
Paper & Date & Pages\\[0.35em]
\midrule
% CONTENTS_ROWS_1
\bottomrule
\end{tabular}
\endgroup
\vfill
{\small Continued on the following page.}

\newpage
\fronttitle{Contents, continued}
{\small Optics, radiometry and astronomical models\par}
\vspace{1.0em}
\begingroup
\small
\renewcommand{\arraystretch}{1.2}
\begin{tabular}{@{}p{0.68\textwidth}@{\hspace{0.03\textwidth}}p{0.18\textwidth}@{\hspace{0.03\textwidth}}>{\raggedleft\arraybackslash}p{0.08\textwidth}@{}}
\toprule
Paper & Date & Pages\\[0.35em]
\midrule
% CONTENTS_ROWS_2
\bottomrule
\end{tabular}
\endgroup
\end{document}
